\section{Bitmap 分配器：\texttt{pmm\_bitmap.c}}

\subsection{核心数据结构}

\codefile{src/kernel/mm/pmm\_bitmap.c}
\begin{lstlisting}[style=cstyle]
struct pmm_region {
    uint32_t base;               // 区域起始物理地址
    uint32_t pages;              // 管理的总页数
    uint8_t* bitmap;             // 位图指针（每 bit 对应一页）
    uint32_t bitmap_bytes;       // 位图有效字节数 = ceil(pages/8)
    uint32_t bitmap_storage_bytes; // 实际分配字节数（对齐后）
    uint32_t bitmap_pages;       // 位图本身占用的页数
    uint32_t free;               // 当前空闲页数
};

static struct pmm_region g_regions[PMM_MAX_REGIONS];  // 最多 32 个区域
static uint32_t g_region_count = 0;
\end{lstlisting}

\begin{figure}[H]
  \centering
  % figures/ch07/fig03-bitmap.tex — auto-extracted from chapter source
\begin{tikzpicture}[node distance=0cm]
  % 内存区域
  \fill[green!20] (0, 0) rectangle (10, 1);
  \node[font=\small] at (5, 0.5) {可用物理内存区域};

  % bitmap 在区域开头
  \fill[blue!30] (0, 0) rectangle (0.8, 1);
  \node[font=\tiny, rotate=90] at (0.4, 0.5) {bitmap};

  % 管理的页
  \foreach \i in {1,...,9} {
    \draw[thin, gray] (\i, 0) -- (\i, 1);
  }

  % 标注
  \draw[decorate, decoration={brace, amplitude=5pt, mirror}]
    (0, -0.1) -- (0.8, -0.1) node[midway, below=5pt, font=\scriptsize] {bitmap\_pages};
  \draw[decorate, decoration={brace, amplitude=5pt, mirror}]
    (0.8, -0.1) -- (10, -0.1) node[midway, below=5pt, font=\scriptsize] {可分配的物理页};

  \node[font=\tiny, above] at (0, 1.1) {\texttt{base}};
  \node[font=\tiny, above] at (10, 1.1) {\texttt{base + pages * 4K}};

  % 位图 bit 示意
  \node[font=\tiny\ttfamily, above=0.3cm, text=blue!60!black] at (0.4, 1.1) {每 bit 对应一页};
\end{tikzpicture}

  \caption{Bitmap 位图放置在区域开头，自身占用几页物理内存}
\end{figure}

\subsection{页 → bit 的映射}

每个物理页对应位图中的 1 bit：
\begin{itemize}
  \item bit = 1：页被占用（used/reserved）
  \item bit = 0：页空闲（free）
\end{itemize}

\codefile{src/kernel/mm/pmm\_bitmap.c}
\begin{lstlisting}[style=cstyle]
// 测试第 bit 位是否为 1
static int bitmap_test_r(const struct pmm_region* r, uint32_t bit) {
    uint32_t byte_idx = bit / 8;
    uint8_t  bit_mask = (uint8_t)(1u << (bit % 8));
    return (r->bitmap[byte_idx] & bit_mask) != 0;
}

// 置位（标记为 used）
static void bitmap_set_r(struct pmm_region* r, uint32_t bit) {
    r->bitmap[bit / 8] |= (uint8_t)(1u << (bit % 8));
}

// 清位（标记为 free）
static void bitmap_clear_r(struct pmm_region* r, uint32_t bit) {
    r->bitmap[bit / 8] &= (uint8_t)~(1u << (bit % 8));
}
\end{lstlisting}

\begin{figure}[H]
  \centering
  % figures/ch07/fig04-8-bit-8.tex — auto-extracted from chapter source
\begin{tikzpicture}[node distance=0cm]
  % 位图的一个字节
  \foreach \i/\val in {0/1, 1/1, 2/0, 3/0, 4/1, 5/0, 6/0, 7/0} {
    \node[bitcell, fill=white, minimum width=0.9cm] (b\i) at (\i*1.0, 0) {\val};
    \node[font=\tiny\ttfamily, above=0.1cm of b\i] {bit \i};
  }
  \node[font=\small\ttfamily, left=0.5cm of b0] {bitmap[k]:};

  % 对应的物理页
  \node[memcell, minimum width=0.9cm, minimum height=0.5cm, fill=red!15, font=\tiny] (p0) at (0, -1.2) {used};
  \node[memcell, minimum width=0.9cm, minimum height=0.5cm, fill=red!15, font=\tiny] (p1) at (1.0, -1.2) {used};
  \node[memcell, minimum width=0.9cm, minimum height=0.5cm, fill=green!15, font=\tiny] (p2) at (2.0, -1.2) {free};
  \node[memcell, minimum width=0.9cm, minimum height=0.5cm, fill=green!15, font=\tiny] (p3) at (3.0, -1.2) {free};
  \node[memcell, minimum width=0.9cm, minimum height=0.5cm, fill=red!15, font=\tiny] (p4) at (4.0, -1.2) {used};
  \node[memcell, minimum width=0.9cm, minimum height=0.5cm, fill=green!15, font=\tiny] (p5) at (5.0, -1.2) {free};
  \node[memcell, minimum width=0.9cm, minimum height=0.5cm, fill=green!15, font=\tiny] (p6) at (6.0, -1.2) {free};
  \node[memcell, minimum width=0.9cm, minimum height=0.5cm, fill=green!15, font=\tiny] (p7) at (7.0, -1.2) {free};
  \node[font=\small, left=0.5cm of p0] {物理页:};

  % 箭头
  \foreach \i in {0,...,7} {
    \draw[gray, thin, -{Stealth[length=2pt]}] (b\i.south) -- (p\i.north);
  }
\end{tikzpicture}

  \caption{位图的一个字节（8 bit）对应 8 个连续的物理页}
\end{figure}

\subsection{分配算法：Round-Robin + 线性扫描}

\codefile{src/kernel/mm/pmm\_bitmap.c}
\begin{lstlisting}[style=cstyle]
void* pmm_bitmap_alloc_page(void) {
    // Round-robin: 从上次分配的 region 开始，避免总是扫描第一个 region
    for (uint32_t tries = 0; tries < g_region_count; tries++) {
        uint32_t ri = (g_region_rr + tries) % g_region_count;
        struct pmm_region* r = &g_regions[ri];
        if (r->free == 0) continue;

        uint32_t idx;
        if (pmm_find_free_index(r, &idx)) {
            bitmap_set_r(r, idx);       // 标记为 used
            r->free--;
            g_pmm_total_free--;
            g_region_rr = ri;           // 记住本次从哪个 region 分配的
            return (void*)(uintptr_t)region_index_to_addr(r, idx);
        }
    }
    return NULL;  // OOM
}
\end{lstlisting}

\texttt{pmm\_find\_free\_index} 的核心：先按字节扫描（跳过全 0xFF 的字节），
找到有空位的字节后再逐 bit 检查。这把线性扫描的常数因子降低了 8 倍。

\begin{thinkbox}
Round-robin 跨 region 分配有什么好处？如果总是从 region 0 开始扫描，
region 0 会很快碎片化，而 region 1 始终空闲——分配不均匀。
Round-robin 均衡了各 region 的使用率。
\end{thinkbox}

\subsection{释放算法}

\codefile{src/kernel/mm/pmm\_bitmap.c}
\begin{lstlisting}[style=cstyle]
void pmm_bitmap_free_page(void* page) {
    uint32_t addr = (uint32_t)(uintptr_t)page;
    // 安全检查：对齐？在管理范围内？
    if (!pmm_addr_is_page_aligned(addr)) return;

    // 找到包含该地址的 region
    for (uint32_t i = 0; i < g_region_count; i++) {
        struct pmm_region* r = &g_regions[i];
        if (!region_contains_addr(r, addr)) continue;

        uint32_t idx = region_addr_to_index(r, addr);
        if (!bitmap_test_r(r, idx)) {
            printk("[pmm] double free: 0x%08x\n", addr);  // 防御
            return;
        }
        bitmap_clear_r(r, idx);
        r->free++;
        g_pmm_total_free++;
        return;
    }
}
\end{lstlisting}

\begin{warnbox}
Double-free 检测：如果 bit 已经是 0 还要清，说明同一个页被 free 了两次。
这是严重 bug，打印警告但不执行清除（避免 free count 错乱）。
\end{warnbox}

\subsection{初始化流程}

\texttt{pmm\_bitmap\_init()} 的完整步骤：

\begin{enumerate}
  \item 解析 Multiboot2 mmap tag，找出所有 type=1（Available）的区域
  \item 对每个区域：截断到 $\ge$ 1MiB、$\le$ 4GiB、页对齐
  \item 计算位图大小和放置位置（\texttt{pmm\_compute\_region\_layout}）——
    位图本身也要占物理内存，需要避开内核镜像和 mb2 info
  \item 清零位图 → 将位图自身占用的页标记为 used
  \item 保留内核镜像区域（\texttt{\_\_kernel\_phys\_start} 到 \texttt{\_\_kernel\_phys\_end}）
  \item 保留 Multiboot2 info 区域
  \item 重新计算 free count
  \item Selftest：分配 4 页 → 写入 pattern → 读回验证 → 释放 → 检查 free 恢复
\end{enumerate}

\begin{designbox}
位图放置是初始化中最复杂的部分。位图需要放在可用物理内存中，
但它不能覆盖内核镜像或 Multiboot2 info。
\texttt{pmm\_compute\_region\_layout} 最多尝试 4 次不同的放置方案，
直到找到一个不和任何已有数据重叠的位置。
\end{designbox}

% ============================================================================

